\documentclass[11pt,a4paper]{article}
\usepackage{amsmath,amssymb,amsthm}
\usepackage[margin=1in]{geometry}
\usepackage{hyperref}

\newtheorem{theorem}{Theorem}
\newtheorem{lemma}{Lemma}
\newtheorem{proposition}{Proposition}
\newtheorem{corollary}{Corollary}
\newtheorem{remark}{Remark}

\DeclareMathOperator{\re}{Re}

\title{Analytical Notes: Collision Energy Monotonicity\\
and a Characterization of $\Lambda$}
\author{Working notes --- Session 16}
\date{March 2026}

\begin{document}
\maketitle

\section{Setup}

Zeros $z_1 < z_2 < \cdots < z_N$ evolve under the Coulomb ODE
\begin{equation}\label{eq:ode}
  \dot{z}_j = 2\sum_{k \neq j} \frac{1}{z_j - z_k}, \qquad j = 1,\ldots,N.
\end{equation}
Define the \emph{Coulomb energy}
\begin{equation}\label{eq:energy}
  E(t) = -\sum_{j < k} \log(z_j(t) - z_k(t))^2
       = -2\sum_{j < k} \log|z_j - z_k|.
\end{equation}

\begin{lemma}[Energy identity]\label{lem:energy}
  Under the ODE \eqref{eq:ode},
  \begin{equation}\label{eq:dEdt}
    \frac{dE}{dt} = -4\sum_{j < k} \frac{1}{(z_j - z_k)^2}.
  \end{equation}
\end{lemma}

\begin{proof}
Compute:
\begin{align*}
  \frac{dE}{dt}
  &= -2\sum_{j<k} \frac{\dot{z}_k - \dot{z}_j}{z_k - z_j}.
\end{align*}
Now
\[
  \dot{z}_k - \dot{z}_j
  = 2\sum_{l \neq k}\frac{1}{z_k - z_l}
  - 2\sum_{l \neq j}\frac{1}{z_j - z_l}.
\]
Isolate the mutual term ($l = j$ in the first sum, $l = k$ in
the second):
\[
  = \frac{2}{z_k - z_j} + \frac{2}{z_j - z_k}
  + 2\sum_{l \neq j,k}\left(\frac{1}{z_k - z_l} - \frac{1}{z_j - z_l}\right)
  = \frac{4}{z_k - z_j}
  + 2\sum_{l \neq j,k}\frac{z_j - z_k}{(z_k - z_l)(z_j - z_l)}.
\]
Wait --- the mutual terms: in $\dot{z}_k$, the $l=j$ term gives
$2/(z_k - z_j)$. In $\dot{z}_j$, the $l=k$ term gives
$2/(z_j - z_k) = -2/(z_k - z_j)$. So the difference is
$2/(z_k-z_j) - (-2/(z_k-z_j)) = 4/(z_k-z_j)$.

For the remaining terms ($l \neq j,k$):
\[
  \frac{1}{z_k - z_l} - \frac{1}{z_j - z_l}
  = \frac{z_j - z_k}{(z_k - z_l)(z_j - z_l)}.
\]
Therefore
\[
  \frac{\dot{z}_k - \dot{z}_j}{z_k - z_j}
  = \frac{4}{(z_k - z_j)^2}
  + \frac{2}{z_k - z_j}\sum_{l \neq j,k}
    \frac{z_j - z_k}{(z_k - z_l)(z_j - z_l)}
  = \frac{4}{(z_k - z_j)^2}
  - 2\sum_{l \neq j,k}\frac{1}{(z_k - z_l)(z_j - z_l)}.
\]
Summing over $j < k$:
\[
  \frac{dE}{dt} = -2\sum_{j<k}\frac{4}{(z_k-z_j)^2}
  + 4\sum_{j<k}\sum_{l\neq j,k}\frac{1}{(z_k-z_l)(z_j-z_l)}.
\]
The triple sum vanishes by symmetry. Relabel: for each
ordered triple $(j,k,l)$ with $j < k$ and $l \neq j,k$,
the term is $1/[(z_k-z_l)(z_j-z_l)]$. Consider the
permutation $(j,k,l) \to (l,k,j)$ (swapping the roles of
$j$ and $l$). If $l < k$, this contributes
$1/[(z_k-z_j)(z_l-z_j)]$. The sum over all ordered triples
involves terms of the form
\[
  T_{jkl} = \frac{1}{(z_k-z_l)(z_j-z_l)}
\]
for distinct $j,k,l$. By the partial fraction identity
\[
  \frac{1}{(a-c)(b-c)} + \frac{1}{(b-a)(c-a)} + \frac{1}{(c-b)(a-b)} = 0,
\]
the triple sum vanishes when summed over all cyclic
permutations. Since our sum ranges over all $(j,k)$ pairs
with $j < k$ and all $l \neq j,k$, it equals
$\sum_{\{j,k,l\}} [T_{jkl} + T_{klj} + T_{ljk}] = 0$
(each unordered triple contributes the three cyclic terms,
which sum to zero).

Therefore $dE/dt = -4\sum_{j<k} (z_k - z_j)^{-2}$.
\end{proof}

\begin{corollary}\label{cor:monotone}
  $E(t)$ is strictly decreasing in $t$ as long as all zeros
  are distinct. The Coulomb energy decreases under forward
  heat flow (zeros spread apart, log-distances increase).
\end{corollary}

\section{Energy at $t = 0$ and the zero-free region}

Define the \emph{Coulomb self-energy} of the zeta zeros:
\[
  E_N = -2\sum_{1 \leq j < k \leq N} \log|\gamma_j - \gamma_k|,
\]
where $\gamma_1 < \gamma_2 < \cdots$ are the positive imaginary
parts of the zeta zeros. By Lemma~\ref{lem:energy}, $E_N(t)$
is decreasing in $t$, so
\begin{equation}\label{eq:E-bound}
  E_N(0) > E_N(t) \quad \text{for all } t > 0.
\end{equation}

\begin{proposition}[Energy characterization of $\Lambda$]
\label{prop:energy-lambda}
  The de Bruijn--Newman constant satisfies
  \[
    \Lambda = \inf\{t : E_N(t) \text{ is finite for all } N\}.
  \]
  Equivalently, $\Lambda$ is the supremum of $t$ values at which
  $E_N(t) \to -\infty$ (collision: two zeros merge, creating a
  logarithmic singularity).
\end{proposition}

\begin{proof}
  $E_N(t) = -\infty$ iff $z_j(t) = z_k(t)$ for some $j \neq k$,
  i.e., two zeros have collided. By definition, $\Lambda$ is the
  supremum of $t$ at which $H_t$ has a non-real zero, which occurs
  precisely at collision points.
\end{proof}

\section{The energy rate and a sum rule}

From \eqref{eq:dEdt}, the energy dissipation rate is
\[
  -\frac{dE}{dt}\bigg|_{t=0}
  = 4\sum_{j<k} \frac{1}{(\gamma_j - \gamma_k)^2}
  = 4 \cdot \frac{1}{2}\sum_{j \neq k}\frac{1}{(\gamma_j-\gamma_k)^2}.
\]
The sum $\sum_{j \neq k} (\gamma_j - \gamma_k)^{-2}$ is directly
related to the second moment of the pair correlation:
\begin{equation}\label{eq:sum-rule}
  \sum_{\substack{j,k \leq N \\ j \neq k}}
  \frac{1}{(\gamma_j - \gamma_k)^2}
  = \frac{N}{d^2}\int_0^\infty \frac{R_2(r)}{r^2}\,dr + O(N/d^2)
  \approx \frac{N \cdot C_{\mathrm{GUE}}}{d^2},
\end{equation}
where $C_{\mathrm{GUE}} = \int_0^\infty R_2(u)/u^2\,du \approx 3.255$.

This connects the energy dissipation rate at $t = 0$ to our
collision correction constant! The initial rate at which the
Coulomb energy decreases under forward heat flow is governed
by the same GUE integral that controls the mean-field collision
correction.

\begin{theorem}[Energy--collision duality]\label{thm:duality}
  The mean-field collision correction factor $\beta = 4C_{\mathrm{GUE}}/d^2$
  and the Coulomb energy dissipation rate are related by
  \begin{equation}
    -\frac{dE_N}{dt}\bigg|_{t=0} = 2N\beta + O(N/d^2).
  \end{equation}
  In particular, the collision dynamics (Section 3 of the main paper)
  and the energy monotonicity (this section) are dual descriptions
  of the same GUE structure.
\end{theorem}

\section{Toward an unconditional result}

The energy identity $dE/dt = -4\sum (z_j - z_k)^{-2}$ is
\emph{unconditional} --- it holds for any configuration of points
evolving under the Coulomb ODE. It does not assume RH, Montgomery,
or GUE.

\begin{proposition}[Unconditional energy bound]
  \label{prop:unconditional}
  For the first $N$ zeta zeros evolving under the Coulomb ODE,
  define
  \[
    S_N = \sum_{\substack{j,k \leq N \\ j \neq k}}
    \frac{1}{(\gamma_j - \gamma_k)^2}.
  \]
  Then the collision time for the closest pair satisfies
  \begin{equation}\label{eq:tcoll-energy}
    |t_{\mathrm{coll}}| \leq \frac{E_N(0) - E_N^{\min}}{4S_N},
  \end{equation}
  where $E_N^{\min}$ is the minimum energy achievable (the
  equilibrium energy of $N$ points under logarithmic repulsion
  on the relevant interval).
\end{proposition}

\begin{proof}
  By the energy identity, $E$ decreases at rate at least $4S_N$
  per unit $|t|$ going backward. The collision occurs when $E$
  diverges to $-\infty$, which requires at least
  $(E_N(0) - E_N^{\min})/4S_N$ units of $|t|$.

  (This is crude --- $S_N$ changes as the zeros move, so the
  actual bound involves integrating the time-dependent rate.
  The inequality uses the \emph{initial} rate as a lower bound
  on the average rate.)
\end{proof}

\begin{remark}
  Both $E_N(0)$ and $S_N$ are computable from the known zeta
  zeros without any unproven conjecture. This gives an
  \emph{unconditional} upper bound on the collision time, and
  hence a lower bound on $\Lambda$:
  \[
    \Lambda \geq -\frac{E_N(0) - E_N^{\min}}{4S_N}.
  \]
  The question is whether this bound converges to zero as
  $N \to \infty$. If it does, combined with $\Lambda \geq 0$,
  it gives no new information. But if the bound can be made
  \emph{positive} for some $N$, it would prove $\Lambda > 0$
  --- which contradicts $\Lambda \geq 0$ unless $\Lambda = 0$,
  i.e., RH.

  (Of course, making the bound positive from below is
  extremely unlikely with such a crude estimate.)
\end{remark}

\section{The triple-sum vanishing and its consequences}

The key algebraic fact in Lemma~\ref{lem:energy} is that the
triple interaction sum vanishes identically:
\[
  \sum_{j<k}\sum_{l \neq j,k}
  \frac{1}{(z_k - z_l)(z_j - z_l)} = 0.
\]
This is a consequence of the partial fraction identity
\[
  \frac{1}{(a-c)(b-c)} + \frac{1}{(b-a)(c-a)}
  + \frac{1}{(c-b)(a-b)} = 0,
\]
which holds for any distinct $a, b, c$.

This vanishing is remarkable: it means the Coulomb energy is
an \emph{exact} Lyapunov function for the ODE \eqref{eq:ode},
with dissipation rate given purely by the pairwise sum of
inverse-square distances. There are no three-body corrections.

\begin{corollary}[Exact gap equation]
  The gap $g_{jk}(t) = z_k(t) - z_j(t)$ for any pair satisfies
  \begin{equation}\label{eq:gap-exact}
    \frac{dg_{jk}}{dt} = \frac{4}{g_{jk}}
    + 2\sum_{l \neq j,k}\frac{g_{jk}}{(z_k - z_l)(z_j - z_l)}.
  \end{equation}
  The correction term (the sum) does \emph{not} vanish for
  individual pairs, only when summed over all pairs weighted
  by $1/g_{jk}$. This is why the two-body approximation is
  excellent for the collision time ($\sim$5\% error) but the
  energy identity is exact.
\end{corollary}

\section{Connection to the xi function: an exact identity}

The energy dissipation rate at $t = 0$ can be expressed in
terms of the xi function without reference to individual zeros.

\begin{proposition}
  For $H_0(z) = \xi(1/2+iz)$, the Coulomb energy rate satisfies
  \begin{equation}\label{eq:xi-identity}
    -\frac{dE}{dt}\bigg|_{t=0}
    = 4\sum_{j<k}\frac{1}{(\gamma_j - \gamma_k)^2}
    = 2\int_{-\infty}^{\infty} \left[-f''(z)\right] dz
    - 2\sum_j \lim_{\epsilon \to 0}
      \left[\frac{1}{\epsilon^2} + f''(\gamma_j + \epsilon)\right],
  \end{equation}
  where $f(z) = \log|H_0(z)|$ and the integral and limit extract
  the regular and singular parts of $f''$ respectively.
\end{proposition}

\begin{proof}[Proof sketch]
  From the Hadamard product (assuming real zeros):
  \[
    f''(z) = -\sum_k \frac{2(\gamma_k^2 + z^2)}{(\gamma_k^2 - z^2)^2}.
  \]
  Near $z = \gamma_j$: $f''(z) \approx -1/(z-\gamma_j)^2
  + \sum_{k \neq j} [\text{regular}]$. The regular part at
  $z = \gamma_j$ involves $\sum_{k \neq j} 1/(\gamma_j-\gamma_k)^2$,
  which is exactly our energy dissipation integrand.
\end{proof}

\begin{remark}
  The quantity $-f''(z)$ is our log-concavity measure. This
  identity shows that the \emph{integral} of the log-concavity
  over the real line equals the energy dissipation rate (up to
  the pole subtraction). Our numerical finding that $f''(z) < 0$
  everywhere means every point contributes positively to the
  energy dissipation. The identity makes this connection exact.
\end{remark}

\end{document}
